\section{时空、物质与参考系}
\label{r2ch01:sec:spacetime}

运动学首先回答“物体在何处、怎样移动”，暂不追问运动为何发生。为了使“位置”和“时刻”成为可测量的量，必须先规定空间、时间、物质对象和参考系。经典力学把空间和时间看作物质运动的舞台；相对论则进一步指出，时空结构本身与物质、能量不可完全割裂。二轮学习仍以经典运动学为主，但应清楚经典模型的适用边界。

从伽利略把可重复实验引入物理学，到牛顿力学、分析力学、经典场论和相对论的建立，“物理过程发生在怎样的时空中”始终是理论的第一层结构。经典理论常把时空与物质分开处理：时空提供事件的位置和间隔，物质则在每一个时空点上具有质量、密度、速度、场强等状态。数学上，连续、光滑并带有度量的时空可用流形描述；在本章中不需要流形语言，但要保留一个核心观念：坐标只是给时空点编号的工具，时空点之间可测的距离或间隔才具有直接的物理意义。

时空（spacetime）与物质（matter）的这种区分是一种组织理论的方法，并不要求时空永远是没有动力学的固定背景。
例如广义相对论把引力和时空几何联系起来，有些理论还探讨额外空间维度；
这些研究并不改变本章采用“三维空间加一维时间”的近似。
数学上，带非退化洛伦兹号差度量的相对论时空是一类伪黎曼流形（pseudo-Riemannian manifold）；
经典绝对时间与空间度量则是不同的几何结构，不能只因它也有四个坐标，就认定它与相对论具有同一种度量。

所谓度量，是一条把相邻点坐标差转化为实际间隔的规则。选用不同坐标不会改变实际间隔，却会改变度量公式的写法。因此“建立坐标系”不是改变空间，而是在同一空间上选择一套描述。直角坐标之所以方便，正因为它把平直空间的距离写成最简单的平方和；曲线坐标则牺牲这种形式，换取与问题对称性或几何边界更贴合的坐标线。
这里的度量对应英文 metric；三维空间的直角坐标也称笛卡尔坐标（Cartesian coordinates）。

\subsection{空间与时间}
\label{r2ch01:sec:space-time-coordinates}

先把坐标值与被标记的对象区分开。在选定长度、时间单位后，三维空间点、时刻和事件分别用
\[
  (x,y,z)\in\mathbb R^3,\qquad t\in\mathbb R,
  \qquad (x,y,z,t)\in\mathbb R^4
\]
表示；这里 $\mathbb R$ 是实数集合，$\mathbb R^n$ 是由 $n$ 个实数组成的有序数组集合。
写成四个实数不等于四个坐标具有相同量纲，更不自动规定它们之间应使用哪一种距离公式。
例如用面前、右方和头顶确定三个互相垂直的方向，再选定一把尺，才能把几何位置转换为有单位的数值。

经典力学中的空间具有三维、连续、均匀、各向同性和欧几里得等性质。若选定原点和三个相互垂直的方向，一个空间点可用直角坐标 $(x,y,z)$ 表示。对于具有轴对称或球对称的系统，柱坐标 $(\rho,\varphi,z)$ 或球坐标 $(r,\theta,\varphi)$ 往往更简洁：
\begin{equation}
  x=\rho\cos\varphi,\quad y=\rho\sin\varphi,
  \qquad
  x=r\sin\theta\cos\varphi,\quad
  y=r\sin\theta\sin\varphi,\quad
  z=r\cos\theta .
  \label{r2ch01:eq:001}
\end{equation}
其中 $\theta$ 是从正 $z$ 轴量起的极角，$\varphi$ 是在 $xy$ 平面内从正 $x$ 轴朝正 $y$ 轴量起的方位角。仅与 $\varphi$ 无关表示绕 $z$ 轴的轴对称；同时与 $\theta,\varphi$ 无关才是球对称。三种坐标的几何定义见图~\ref{r2ch01:fig:001}：$OP=r$，$OH=\rho=r\sin\theta$，$HP=z=r\cos\theta$。图中采用三维右手系的平面投影，因此投影后两轴的夹角不一定是直角；$\varphi$ 的角弧仍必须位于 $xy$ 平面，$\theta$ 的角弧必须位于 $OP$ 与 $z$ 轴张成的平面。
柱坐标、球坐标分别对应 cylindrical coordinates、spherical coordinates；
极角和方位角分别对应 polar angle、azimuth angle，方位角也称幅角。
用距离和两个观测方向角确定天体位置，是球坐标在天文观测中的直观应用。

\figCoordinates

时间用来标定事件的先后和间隔。经典力学认为，不同观察者使用同一个绝对时间，因而 $t'=t+t_0$ 只可能相差计时零点。国际单位制中的秒由铯-133原子基态超精细跃迁频率的固定数值定义；米、千克等基本单位也已用固定的基本常数定义。实验室中的尺和钟只是这些定义的实现装置，而不是单位本身。

设惯性系 $S'$ 相对 $S$ 以恒定速度 $\vect V$ 平动，两系坐标轴平行。由
$\vect r'=\vect r-\vect Vt-\vect r_0$，对任意两个事件作差，得到伽利略变换下的关系
\begin{equation}
  \Delta\vect r'=\Delta\vect r-\vect V\Delta t,
  \qquad \Delta t'=\Delta t.
  \label{r2ch01:eq:002}
\end{equation}
只有两个事件同时发生，即 $\Delta t=0$，才有
$|\Delta\vect r'|=|\Delta\vect r|$。因此经典力学保持不变的是时间间隔和同时测量的空间距离，
不是任意两个事件的空间坐标差。例如车厢内一个静止标记在前后两个时刻满足
$\Delta\vect r'=0$，地面观察者却测得 $\Delta\vect r=\vect V\Delta t$；两者并不矛盾。
即使两系坐标轴还相差一个固定转动，同时距离的大小仍不变，因为正交转动不改变矢量长度。
当速度接近光速时，空间和时间必须合并为闵可夫斯基时空。采用度规号差 $(-,+,+,+)$，惯性系之间保持不变的是
\begin{equation}
  \Delta s^2=\lvert\Delta\vect r\rvert^2-c^2\Delta t^2.
  \label{r2ch01:eq:003}
\end{equation}
式~\eqref{r2ch01:eq:003} 是狭义相对论的时空间隔。惯性系之间的洛伦兹变换在适当的低速近似下
回到式~\eqref{r2ch01:eq:002} 的伽利略变换；不能把“时空间隔不变”直接改写成
“不同时刻的空间距离也不变”。

直角坐标中两个点 $A(x_1,y_1,z_1)$、$B(x_2,y_2,z_2)$ 的距离满足
\[
  AB^2=(x_1-x_2)^2+(y_1-y_2)^2+(z_1-z_2)^2,
\]
微元形式为 $\dd l^2=\dd x^2+\dd y^2+\dd z^2$。坐标原点可以放在任意位置，坐标轴也可作任意刚性转动；这些操作改变 $(x,y,z)$，却不改变 $AB$。这一区分解释了“空间点可以脱离坐标值存在”的含义。

柱坐标中，$\rho$ 是点到 $z$ 轴的距离；球坐标中，$r$ 是点到原点的距离。角坐标具有周期性：$\varphi$ 与 $\varphi+2n\pi$ 表示同一几何方向，但在描述连续转动时不应强行把 $\varphi$ 截回某一个长度为 $2\pi$ 的区间，否则质点每绕轴一周，坐标会发生人为跳变。通常取 $0\leq\theta\leq\pi$，在 $\theta=0,\pi$ 两个极点处，$\varphi$ 的不同数值表示同一点。这不是空间的奇异，而是球坐标本身的坐标奇异。

对称性的判断也必须和变换对应。若标量函数在绕 $z$ 轴任意转角下不变，即不显含 $\varphi$，称为轴对称；若还不依赖极角 $\theta$，只依赖 $r$，才称为球对称或各向同性。柱坐标中的“与 $z$ 无关”表示沿轴平移对称，并不等同于绕轴对称。中心反演对称只要求
\[
  f(x,y,z)=f(-x,-y,-z),
\]
它比球对称弱得多。

绝对时空观把任意时刻看作一个完整的三维空间切片。不同切片按同一时间参数线性排列，理想的空间不随时间整体膨胀、收缩或自发转动。更严格地说，这种结构可用牛顿—嘉当几何（Newton--Cartan geometry）刻画：四维时空上有读取时间间隔的时间一形式，也有只描述空间方向的退化逆度规。
记时间一形式为 $c_\nu$，空间逆度规为 $s^{\mu\nu}$，其基本兼容条件写成
\[
  \sum_{\nu=0}^{3}s^{\mu\nu}c_\nu=0\qquad(\mu=0,1,2,3).
\]
这个式子表达“纯空间的测量结构不包含时间方向”。为看清含义，取 $x^0=t$，在惯性直角坐标中可用
\[
  (c_0,c_1,c_2,c_3)=(1,0,0,0),\qquad
  (s^{\mu\nu})=\operatorname{diag}(0,1,1,1)
\]
表示它们，此时矩阵与列向量的乘积确实为零。这里的 $c_\nu$ 不是光速 $c$，上下标只是标记四维分量。
完整的牛顿—嘉当理论还需要规定比较邻近点方向的联络；上式只介绍其时空测量结构。
中学竞赛不需要使用这些张量，但这个背景能说明为什么经典力学可以给所有观察者规定同一个绝对时间，却不能为不同时刻的两个点定义一个观察者无关的空间距离。
经典理论也可以作协变的几何表述；不能把“不是洛伦兹时空”理解成“没有任何协变性”。

在三维空间坐标作刚性转动时，质量、能量等标量保持数值不变，动量、角动量等矢量按坐标轴一起转动。更一般的伽利略变换还混合位置与时间，速度因此改变而加速度不变。相对论则取消绝对同时性，用光速把时间量纲换算成长度量纲，并以时空间隔区分时间样、光样和空间样事件。事件先后只有在因果允许的范围内才对所有惯性观察者一致；这一点将在相对论章节详细处理。

“数值加单位”共同构成一个物理量的表示。例如两点坐标差给出数值 $(1,1,1)$，若坐标单位为米，距离才是 $\sqrt3\ \mathrm m$，而不是无单位的 $\sqrt3$。米的现代定义把真空光速固定为精确值
$c=299\,792\,458\ \mathrm{m\,s^{-1}}$；秒的定义把铯-133原子基态超精细跃迁频率固定为
$9\,192\,631\,770\ \mathrm{Hz}$。这些定义彼此关联：必须先有可实现的秒，才能用光在一定时间内传播的距离实现米。它们同时依赖目前经受实验检验的相对论和量子理论。

绝对时空中，任意两个事件都有绝对时间差 $|t_2-t_1|$；空间距离只在“同时”的三维切片上直接定义。一个在不同时刻移动的物体所经历的“路长”，必须先说明用哪个参考系的空间切片连接相邻位置。这个事实已经预示：空间坐标虽然看似直观，也始终带有参考系选择。

经典空间允许任意平移和转动而保持欧氏距离不变。把一套直角基矢同时转动，不会改变点积与叉积的几何意义；因此标量关系、矢量关系可以写成与具体轴名无关的形式。所谓物理规律的几何性，正是指它不应依赖我们把胸前方向叫作 $x$ 还是 $y$。相对运动的伽利略变换进一步改变坐标网格随时间的位置，这时速度不再具有不变数值，必须研究量的变换规律。

相对论的度规也可采用相反号差写成
$c^2\dd t^2-\dd x^2-\dd y^2-\dd z^2$；两种约定给出相同因果分类，只要全书保持一致即可。有限光速使相隔足够远的两个事件可能无法用信号互相影响，此时不同惯性系甚至会对它们的时间先后给出不同判断。经典“先发生必定是原因”的时序直觉因而要被更精确的光锥因果关系取代。

坐标的多值性还需要从“运动”而非“静态图形”理解。例如一个点连续绕 $z$ 轴转过三周，它的方向和初始时相同，但连续选取的方位角应当从 $0$ 增加到 $6\pi$，而不是每到 $2\pi$ 就突然跳回零。反过来，球坐标的极点、柱坐标的轴线都会使坐标标记失去唯一性。这些都是描述工具的特性，不是实际空间出现了断裂或重叠。

判断对称性时也应先说清“在什么变换下不变”。绕轴转动不变、沿轴平移不变、对原点反演不变是三种不同的条件。只说“柱对称”容易把前两者混在一起，因此后文尽量直接写成“轴对称”或“沿轴方向均匀”。一个中心反演对称的函数可以在各个方向上截然不同，所以中心对称绝不能代替球对称。

\subsection{物质、事件与场}
\label{r2ch01:sec:matter}

事件是“在某一时刻、某一位置发生的事情”，可用 $(t,\vect r)$ 标记。物体随时间占据的一串事件在时空中形成世界线。研究对象的内部尺度和形变可以忽略时，可把它理想化为质点；多个质点组成质点系；当微观尺度远小于所研究的尺度时，可使用连续介质模型，以密度、速度、温度等场描述它。

事件除了外在的时空坐标，还具有由物质种类决定的内禀状态。例如经典电磁场在一个事件处可用电场和磁场描述；带电质点除位置、速度和质量外还带有电荷。模型随研究层次而改变：电子在经典问题中常近似为带电质点，在量子理论中则要用波函数、自旋和相应的动力学方程描述。模型的层次不同，并不意味着低层次模型毫无价值；关键是它的空间、时间和能量尺度是否覆盖当前问题。

这里还应区分\emph{内禀参数}与\emph{运动状态}：同一种电子具有确定的质量、电荷和自旋量子数，
但它的动量和能量随状态、参考系而改变。磁矩与自旋相关，却不是把“磁矩”与“自旋”当成同一个物理量。
以电磁相互作用为例，除电场、磁场外，还可引入标势 $\phi$ 和矢势 $\vect A$，满足
\[
  \vect B=\nabla\times\vect A,\qquad
  \vect E=-\nabla\phi-\frac{\partial\vect A}{\partial t}.
\]
符号 $\nabla\phi$ 表示标势的空间梯度，$\nabla\times\vect A$ 表示矢势的旋度；
量子力学（quantum mechanics, QM）中的带电粒子通过这些势与电磁场耦合。
势的表示具有规范自由，不能把同一电磁场的不同规范表示误认为不同的可观测场。
量子场论（quantum field theory, QFT）进一步把场量子化，光子是量子化电磁场的激发。

以电子为例可以更清楚地看到“内禀状态”的层次。只研究它在宏观电磁场中的轨迹时，质量和电荷已足以给出经典运动；研究原子光谱和磁矩时，必须加入自旋以及量子态；在相对论量子力学中，状态由四分量狄拉克旋量及其方程描述；在量子电动力学（quantum electrodynamics, QED）中，还要把电子场与电磁场都量子化。同一物理对象的“属性清单”并不是永远固定的，而由要回答的问题决定。

气泡上升可以用来检验这种区分。若系统只包括气泡中的气体，气体质心向上时其动量也向上；
若系统还包括周围液体，总动量就应写成
\[
  \vect P=\int_{V_{\rm gas}}\rho_{\rm gas}\vect v_{\rm gas}\,\dd V
          +\int_{V_{\rm liquid}}\rho_{\rm liquid}\vect v_{\rm liquid}\,\dd V.
\]
液体各处如何流动、容器怎样约束液体，都会影响第二项，不能仅凭气泡的上升方向确定整个系统的动量。
若把气泡等效成“相对于均匀液体的一块密度亏空”，所记的有效量还包含被置换液体的贡献，
不再只是气泡内气体的真实质量乘速度。这一例子说明建模时必须说清系统和变量，而不是说明动量不随参考系改变。

因此，建模不是把真实物体“变成”质点、刚体或连续介质，而是有意识地保留与观测精度有关的自由度。题目中的“忽略大小”、“不计形变”、“均匀介质”都是模型边界。如果答案恰好依赖被删去的大小、转动或密度分布，就必须升级模型，而不能继续沿用质点近似。
本章所说的经典物理学（classical physics）包括牛顿力学、分析力学、经典电磁学、经典统计力学及狭义相对论；
“经典”在这里主要与量子理论相区别，并不等同于“只讨论低速运动”。

质点（point mass, particle）模型把研究对象的尺度、形状和内部转动暂时忽略，只保留一条世界线（world line）及必要的内禀参数。对每个时刻只给出一个位置，正是运动方程
\[
  \vect r=\vect r(t)=\bigl(x(t),y(t),z(t)\bigr)
\]
的含义。质量表征物体在给定相互作用下改变运动状态的难易程度，也是牛顿力学和引力理论的基本参数。
1889年至2019年5月20日前，千克用存放于法国国际计量局的国际千克原器的质量定义。
它是铂铱合金圆柱，直径与高度约为 $39\ \mathrm{mm}$，体积约 $47\ \mathrm{cm^3}$。
复制体需要与原器比较来传递质量单位，而不是每一个复制体都另行定义一份千克。
现行国际单位制已把普朗克常数的数值固定为
\[
  h=6.62607015\times10^{-34}\ \mathrm{kg\,m^2\,s^{-1}},
\]
从而把千克与秒、米及量子电学测量联系起来。历史实现方式可以改变，质量量纲 $M$ 及物理规律不随之改变。
这一实现链先用秒和米确定时间、长度单位，再用固定的 $h$ 确定质量单位；
不能把普朗克常数的单位误写成能量单位，它的量纲是作用量，即能量乘时间。

质量是经典力学中描述惯性和引力性质的标量。若系统的总质量为 $M$，其质心位置定义为
\begin{equation}
  \vect r_C=\frac{1}{M}\sum_i m_i\vect r_i,
  \qquad
  \text{或}\qquad
  \vect r_C=\frac{1}{M}\int \vect r\,\dd m.
  \label{r2ch01:eq:004}
\end{equation}
电磁场、引力场等也是物质存在和相互作用的形式。场量随时空连续分布，并能够携带能量、动量和角动量。图~\ref{r2ch01:fig:002} 用两种横波结构说明：场并非只是一套计算符号，其扰动能够独立传播。

\figFieldWaves

图~\ref{r2ch01:fig:002} 中两种波的“偏振”不是同一类箭头的不同画法。电磁波的 $\vect E$、$\vect B$ 是矢量场，在真空平面波中它们与波矢 $\vect k$ 两两垂直，并且由右手系确定传播方向。弱引力波则是时空度规的张量扰动，$+$ 与 $\times$ 两个模式描述一圈自由质点的横向距离如何周期改变。两者都只有两个独立横向偏振，但一个源自矢量场，一个源自张量场，不能把引力波简单想成另一种电磁波。

具体地，图~\ref{r2ch01:fig:002}(a) 取传播方向为 $+x$，可写成
\[
 \vect E=E_0\sin(kx-\omega t)\uvec e_z,\qquad
 \vect B=-\frac{E_0}{c}\sin(kx-\omega t)\uvec e_y.
\]
两个场共用同一相位，在每一个波节同时为零；波形反号时，两组箭头同时反向，
所以 $\vect E\times\vect B$ 始终指向 $+x$（波节处为零）。图中的场箭头分别按各自幅值缩放，
不能用纸上的箭头长度直接比较量纲不同的电场与磁场。

图~\ref{r2ch01:fig:002}(b) 进一步比较圆偏振的形式结构。螺旋度 $\lambda$ 描述绕传播轴转动时
偏振态所带的旋转相位；两个圆偏振基矢 $e_+,e_-$ 对应 $\lambda=+1,-1$。
符号 $\otimes$ 表示由两个矢量构造二阶张量：同号组合的螺旋度相加为 $\pm2$，异号组合相加为零。
引力波的物理横向无迹部分只有 $\lambda=\pm2$；这里是偏振张量的数学构造，
并不表示引力波由两束电磁波组成。将这两个圆偏振组合为实的线偏振，便得到 $+$ 与 $\times$。

图~\ref{r2ch01:fig:002}(c) 把这种张量性质转成可观察的几何关系：取一圈起初均匀分布的自由质点，
$+$ 模在两个相互垂直方向上交替伸长和压缩；$\times$ 模的主轴相对它转过 $45^\circ$。
在图示相位约定下，$t=0,T/2,T$ 时回到圆形，$T/4$ 和 $3T/4$ 时两条主轴的伸缩互换。
虚线圆是未扰动位置，实线和点是当时的位置；$T$ 表示一个周期。

以波为例还可看出，场的状态必须包含足以预言后续演化的全部分量，只画出某一时刻的波形并不总是完整初值。对波动方程，通常还要给定场的时间变化率；这和质点的二阶运动方程既需初位置又需初速度完全类似。因此“物质模型”与“状态变量”是连在一起的：先选模型，再决定哪些数据构成完整初态。

物理方程必须满足量纲齐次。例如位置、速度、加速度和质量的量纲依次为
\begin{equation}
  [\vect r]=L,\qquad [\vect v]=LT^{-1},\qquad
  [\vect a]=LT^{-2},\qquad [m]=M.
  \label{r2ch01:eq:005}
\end{equation}
量纲检查不能证明公式一定正确，却能迅速排除大量错误。改变单位制只改变数值表示，不改变同一物理规律。

长度、时间和质量是力学中最基本的三个量纲。相同量纲的物理量才能直接相加减，和、差仍保留原量纲；不同量纲可以相乘除，组合成速度、加速度、动量、能量等导出量。指数函数、对数函数和三角函数的自变量必须是无量纲纯数，因此 $\e^{-t/\tau}$ 合法，而 $\e^{-t}$ 只有在 $t$ 已经用某个时间尺度无量纲化后才有意义。量纲法则还提示我们主动寻找无量纲参数，它们常控制一个问题的相似性、极限行为与近似条件。

质点系（particle system）是质点模型的自然推广。把若干质点放在同一系统中后，相互作用可分为内力与外力；动量、角动量和能量的系统性规律由此出现。统计力学进一步研究含有巨大粒子数的系统，不再逐一追踪每条世界线，而在相空间中考虑大量可能状态的概率分布、系综平均和长时间平均。宏观压强、温度和熵正是这种层次上涌现的量。

连续介质（continuous medium）模型把离散粒子在宏观尺度上平滑化。最简单的连续介质至少需要质量密度和速度：
\[
  \rho=\rho(\vect r,t),\qquad \vect v=\vect v(\vect r,t).
\]
前者是标量场，后者是矢量场；弹性体还需要应变和应力，流体还需要压强、黏性等状态量。材料可以表现为弹性、黏性、塑性或黏弹性，选择何种本构关系取决于时间尺度、形变量和材料性质。连续介质理论的重要意义之一，是它把“每个空间点都有物理量”的场论语言系统化。
近平衡条件下，还可以从热力学与统计物理出发建立输运关系。量子理论中的不确定关系则提醒我们，
不能在任意微小尺度都把粒子当作沿精确世界线运动的经典小球；经典轨迹和宏观连续场分别有自己的适用范围。

场概念的历史大致经历两层理解。第一层把场视为描述粒子分布或超距作用的数学工具；第二层则认识到相互作用以有限速度传播，场可以脱离具体源粒子的即时位置传播，并独立携带能量和动量。于是场不只是方便的记号，也是一种物质形态。源粒子改变周围场，另一个粒子再与当地场相互作用，这种局域图像取代了瞬时超距作用。

按照场值的数学类型，可区分标量场、矢量场、张量场和旋量场。经典电磁场常以电磁四势或电、磁场矢量描述；弱引力扰动具有张量性质。自由场最基本的分析方法是傅里叶分解：把复杂空间分布分解为不同波矢、频率和偏振的平面波叠加。对场而言，比“场的总质量”更自然的局域量是能量密度、动量密度、波矢、角频率和色散关系。图~\ref{r2ch01:fig:002} 中电磁波的场矢量偏振与引力波的 $+$、$\times$ 张量偏振正体现了场类型不同所带来的结构差异。

四类经典模型之间并非彼此排斥，而是按问题尺度嵌套。研究行星公转时，行星半径远小于轨道尺度，可以把整颗行星当作质点；研究行星自转和潮汐时，必须恢复刚体或连续介质结构；研究内部温度和对流时，还要加入热力学状态场。模型删去的自由度必须在当前精度下对答案影响足够小，不能因为“质点更容易算”就忽略题目明确依赖的形状或转动。

质点系模型还强调“系统边界”。对一个孤立质点来说，它所受作用都来自系统之外；把施力者一并纳入系统后，原来的作用便成为内部相互作用。这一重新划分不改变客观力，却改变了动量和能量守恒式中哪些项属于边界交换。当质点数极多时，逐个求解世界线既不现实也没有必要，统计力学便在给定宏观约束下对可能微观状态取概率平均。

连续介质则把足够小的宏观体元看成仍含有大量微粒，从而使 $\rho(\vect r,t)$、$\vect v(\vect r,t)$ 等量成为空间和时间的连续函数。介质内部如何传递应力还需要本构关系：弹性、黏性、塑性和黏弹性是不同的近似。场模型又进一步，不必先想象每个场值都是某群粒子的缩写；在局域相互作用图像中，场的扰动本身就能在源离开之后继续传播并携带能量。

自由线性场以平面波为基本模式，与自由质点以匀速直线运动为基本模式类似。傅里叶分解把复杂场展开为不同波矢、角频率和偏振的平面波；场之间有相互作用时，再研究这些模式如何耦合、散射与交换能量。因此场论中比“某块场的质量”更自然的量，是局域能量密度、动量密度、波矢和色散关系。

质点系中的“内力—外力”区分取决于系统边界。两个粒子之间的作用在把两者都纳入系统时是内力，只取其中一个时便成为外力。改变系统边界不改变客观相互作用，却改变动量、能量方程中哪些项被写成流入或流出。连续介质也遵循同一思想：取微小体元作为系统，通过体力、表面力和通量与外界交换物理量。

统计力学中的系统状态仍可从大量质点的坐标和动量出发，但宏观上不可能也无必要逐一测量。把每一个可能微观状态看作相空间中的代表点，依据给定约束建立概率分布，再对宏观量取统计平均，就得到压强、内能等可测结果。这里提到统计观点，是为了说明“物质在时空点上的结构”可以按研究层次选择；本章运动学仍只需要明确研究对象的位置、速度和约束。

连续介质的场 $\rho(\vect r,t)$、$\vect v(\vect r,t)$ 同时依赖空间和时间。固定空间点观察场量随时间变化，得到局部时间导数；跟随某个介质微团观察，还要加上微团穿过空间梯度带来的变化。这种“固定观察者”与“随体观察者”的区别，和后文绝对导数、相对导数、随体导数的区分具有同一结构。

场从“辅助描述”发展为“独立物质”的关键证据，是扰动具有有限传播速度并在离开源之后仍能携带能量。电荷加速产生的电磁波、天体系统产生的引力波都可传播到远处探测器。图~\ref{r2ch01:fig:002} 不能理解为两条普通介质绳波：电磁波中 $\vect E,\vect B$ 与传播方向互相垂直，引力波则表现为空间横向距离的两种张量形变。图中的箭头和偏振符号对应不同物理对象，不能互换。

傅里叶分解把复杂场写成许多平面波模式，是因为线性自由场的每个平面波模式可独立演化。质点自由运动以匀速直线为基本模式，场的自由运动则以给定波矢、频率和偏振的平面波为基本模式；有相互作用时，再研究这些基本模式如何耦合、散射和交换能量。这个类比解释了为什么波矢与角频率在场论中的地位类似于动量与能量。

更现代的量子场论把粒子看成场的量子激发。这个观点超出本章范围，但能避免把“质点”和“场”误认为永远互不相干的两种实体。经典质点世界线、经典连续场都是特定尺度下非常有效的近似；竞赛解题首先要熟练使用这些近似，同时在题设暗示高速、微观或强场条件时意识到模型可能需要升级。

可以把场的响应类比为一张受扰动的弹性薄膜：自由起伏帮助我们想象场的传播，
外界作用引起形变帮助我们想象相互作用。但这只是类比，不表示真空真是一张二维材料膜。
量子场的真空态也有其量子结构，不能用经典“没有振动，所以什么都没有”的图像代替；
真空态、能量及粒子激发的精确定义依赖所采用的场论。

\subsection{参考系与客观规律}
\label{r2ch01:sec:frame}

描述运动必须说明相对于谁观察。参考系包括一个参考物、与它固连的空间标架以及计时装置。坐标数值依赖参考系，但客观物理事件不依赖坐标选择。一个合格的物理规律应在允许的坐标变换下保持形式不变。

坐标系是一套给时空点赋坐标的规则；参考系则还包含实际或理想的观察者、与观察者固连的标架、尺度和时钟。观察者可以平动、转动甚至加速，每一时刻都可在自身附近建立固有标架。不同参考系会给同一事件不同的坐标，但事件本身、两个物体是否相撞、某一时钟实际走过多少固有时等客观事实不因坐标名称改变。

在经典时空中，一个观察者可把自己的位置当作原点，把胸前、右手和头顶三个方向当作相互正交的坐标轴，再用与自己同步的钟标记时间。这套与观察者一起运动的尺和钟构成他的固有参考系。若观察者转身，同一远处物体的坐标分量会改变；若观察者平动，原点到物体的位矢也会改变。但把所有矢量、导数和初值按同一变换一并改写后，客观规律不应依赖“轴叫什么”或“原点放在哪里”。

所以，参考系变换时要分清三个层次：事件本身不变，它的坐标表示会变；单个物理量的数值或分量可能变，例如速度在不同惯性系中不同；完整的物理方程经相应变换后却应保持允许的形式。把“方程形式不变”误解为“所有物理量数值都不变”，会从一开始就写错相对运动关系。

物理规律通常表现为一组在一定条件下反复成立、且不依赖某些偶然参数的关系。对于力学系统，可以选取一组广义坐标 $q_\alpha$ 和广义速度 $\dot q_\alpha$ 描述状态。经典决定论要求：给定某一时刻全部独立的 $q_\alpha(0),\dot q_\alpha(0)$，运动方程应能确定以后以及在可逆理想模型中的以前状态。典型方程写成
\[
  f_i(\ddot q_\alpha,\dot q_\alpha,q_\alpha,t)=0,
\]
独立方程数与自由度数相符，并且最高为二阶。只给位置而不给速度通常不能唯一确定运动；只写微分方程而漏掉初始条件也同样不完整。

还要补上数学上的适用条件：方程须能在讨论的区域内解出
$\ddot q_\alpha=F_\alpha(q,\dot q,t)$，右端对时间连续，并对状态变量满足保证局部唯一性的正则条件，
例如具有连续且局部有界的一阶偏导数。此时由给定初态可在相邻时间区间内唯一延拓。
仅凭“方程数相等、最高二阶”并不能自动保证唯一性；遇到碰撞、奇点或约束切换，还要补充连接条件。

“知道初态便能预言整个演化”的思想曾被拉普拉斯用“全知者”形象化：若能精确知道宇宙中所有粒子的状态，原则上便能由规律推出过去和未来。这个观点对经典物理的建模影响深远，也引出了自由意志、耗散、混沌和量子随机性等问题。这里不讨论它的哲学结论，只保留其科学方法上的含义：可检验理论必须给出从已知状态到可观测结果的确定规则或概率规则。

若理想系统的方程不显含绝对时刻，便具有时间平移对称性；若忽略耗散后方程在 $t\mapsto-t$、速度反号下保持形式，便具有时间反演对称性。真实摩擦和不可逆宏观过程会让简单的反演图像失效，但不能因此在运动学层面随意丢掉初始条件或方向信息。
时间反演与无耗散也不能直接画等号。例如固定外磁场中的洛伦兹力与速度成正比且始终不做功，
只反转速度而不反转磁场，一般不能得到同一外场下的反向轨迹；反演时必须同时检查外部条件如何变换。

“方程至多二阶”的表述需要结合状态选取来看。若把位置作为广义坐标，牛顿力学通常给出二阶方程，所以必须同时给定初位置和初速度；若把位置与速度一起当作状态变量，同一系统也可改写为一阶方程组。两种写法的初值数目相同，都不能用一个“初始位置”代替完整状态。方程数目、自由度数目和独立约束数目必须匹配，否则便会出现无限多解或互相矛盾的条件。

决定论并不等于实验上能做到无限精确的预言。混沌系统对初值误差极敏感，耗散环境可使我们不再掌握封闭系统的全部自由度，量子测量还引入概率性。在满足上述正则条件的经典运动学模型中，完整的初值和约束应当确定局部唯一的运动；若答案仍任意，首先要查题设是否漏条件、约束是否退化、方程是否到达奇点，而不是把任意性推给物理规律。

经典力学中的惯性系彼此由伽利略变换联系。若 $S'$ 相对于 $S$ 以恒定速度 $\vect V$ 平动，且两系坐标轴始终平行，则
\begin{equation}
  \vect r'=\vect r-\vect Vt-\vect r_0,\qquad
  \vect v'=\vect v-\vect V,\qquad
  \vect a'=\vect a.
  \label{r2ch01:eq:006}
\end{equation}
因此加速度在伽利略变换下不变。若一个封闭系统的运动方程在所有惯性系中形式相同，就体现了伽利略相对性原理。确定初始位置和初始速度后，二阶运动方程才给出唯一运动；这也是运动学初始条件不能遗漏的原因。

伽利略相对性原理并不是说所有物理量在不同惯性系中数值相同，而是说封闭实验的基本规律具有相同形式。位置和速度随参考系改变，质量、时间间隔和加速度在经典变换下不变；动量和能量的数值也会改变，但相应守恒定律仍保持形式。区分“量的数值不变”和“方程形式不变”，是以后理解协变性与守恒律的基础。

\section{质点运动的描述}
\label{r2ch01:sec:kinematics}

\subsection{位置、速度、加速度与加加速度}

在给定参考系中，质点的位置是时间的矢量函数
\begin{equation}
  \vect r=\vect r(t).
  \label{r2ch01:eq:007}
\end{equation}
轨迹是消去参数 $t$ 后得到的几何曲线，运动方程则保留了质点在曲线上运动的快慢和方向。平均速度与瞬时速度、加速度、加加速度分别为
\begin{equation}
  \bar{\vect v}=\frac{\Delta\vect r}{\Delta t},\qquad
  \vect v=\frac{\dd\vect r}{\dd t},\qquad
  \vect a=\frac{\dd\vect v}{\dd t},\qquad
  \vect j=\frac{\dd\vect a}{\dd t}.
  \label{r2ch01:eq:008}
\end{equation}
速度沿轨迹切线方向，加速度不一定与速度同向。图~\ref{r2ch01:fig:003} 同时画出位置空间与速度空间，后者中曲线的切向量是加速度。

瞬时量必须理解为极限，而不是“把很小的 $\Delta t$ 直接当成零”。更完整地写，
\[
  \vect v(t)=\lim_{\Delta t\to0}
  \frac{\vect r(t+\Delta t)-\vect r(t)}{\Delta t},
  \qquad
  \vect a(t)=\lim_{\Delta t\to0}
  \frac{\vect v(t+\Delta t)-\vect v(t)}{\Delta t}.
\]
第一个极限把位置曲线上的割线方向变成切线方向；第二个极限则在速度空间中做同一件事。若只看速率 $|\vect v|$ 的增减，会遗漏速度方向变化对加速度的贡献；匀速圆周运动就是最简单的反例。

\figKinematicVectors

图~\ref{r2ch01:fig:003} 采用同一个可直接求导的运动，避免把两个空间的曲线误认为任意示意。
取长度尺度 $\ell$、时间尺度 $\tau$，令 $u=t/\tau$，则
\[
 \begin{aligned}
 \frac{\vect r}{\ell}&=\bigl(u+0.1u^2,\ 0.1u^3-0.6u\bigr),\\
 \frac{\vect v}{\ell/\tau}&=\bigl(1+0.2u,\ 0.3u^2-0.6\bigr),\\
 \frac{\vect a}{\ell/\tau^2}&=\bigl(0.2,\ 0.6u\bigr).
 \end{aligned}
\]
点 $P$ 与速度端点 $V$ 都对应 $u=0.8$。左图沿时间增加方向的切线就是 $\vect v$ 的方向；
右图从速度端点 $V$ 出发的切线就是 $\vect a$ 的方向。
这里的 $\vect a$ 不平行于 $\vect v$，因而它既改变速率，又改变运动方向。

矢量函数是从时间轴到二维或三维矢量空间的映射。每一个 $t$ 对应唯一的 $\vect r(t)$；它在位置空间中描出的像集是轨迹。平均速度是有限时间内位移与时间之比，瞬时速度则是割线趋于切线的极限。二阶中心差分
\[
  \vect a(t)=\lim_{\Delta t\to0}
  \frac{\vect r(t+\Delta t)-2\vect r(t)+\vect r(t-\Delta t)}{\Delta t^2}
\]
清楚表明，加速度衡量相邻等时间位移之差，而不是简单等于“单位时间内速度大小的变化”。即使速率不变，速度方向改变仍会产生加速度。

把各时刻的速度矢量移到同一原点，并连接其端点，得到速度空间中的曲线。速度空间中端点的切向变化率就是加速度。在匀强重力场中，抛体的速度端点沿竖直直线匀速移动；在有心力问题中，速度空间也常把位置空间中复杂的轨迹变为更简单的曲线。更一般的动量空间是近代物理和分析力学的重要工具。

一个具体例子是开普勒运动的速度图。设引力加速度为 $\vect a=-\mu\vect r/r^3$，
单位质量的角动量 $h=r^2\dot\theta$ 非零；$\mu$ 为引力参数，不是质量。
在轨道平面内取固定直角轴，由链式法则
\[
 \frac{\dd\vect v}{\dd\theta}
 =\frac{-\mu\uvec e_r/r^2}{h/r^2}
 =-\frac{\mu}{h}(\cos\theta,\sin\theta).
\]
逐分量积分，得到
\[
 \vect v=\vect C+\frac{\mu}{h}(-\sin\theta,\cos\theta),\qquad
 |\vect v-\vect C|=\frac{\mu}{|h|},
\]
其中 $\vect C$ 是由初态决定的恒矢量。因此轨道在速度空间中位于一条圆上，
圆心一般不在速度原点；椭圆轨道完成一周时描出整圆，非束缚轨道通常只对应圆弧。
半径 $\mu/|h|$ 的量纲是速度。它说明位置空间中的非圆轨道也可以有圆形的速度图，
并不是说质点在真实空间做匀速圆周运动。

加加速度 $\vect j$ 描述加速度随时间的变化，在快速机构、舒适性设计和辐射问题中会出现。继续求导还能定义更高阶运动量。高阶导数并不改变“运动由位置函数决定”的基本事实，只是在不同阶次上提取局部变化信息。
加加速度也称急动度，对应英文 jerk；它等于位置对时间的三阶导数：$\vect j=\dd^3\vect r/\dd t^3$。

在直角坐标系中，固定基矢的时间导数为零，所以求导可以逐分量进行：
\begin{equation}
  \vect r=x\uvec e_x+y\uvec e_y+z\uvec e_z,\quad
  \vect v=\dot x\uvec e_x+\dot y\uvec e_y+\dot z\uvec e_z,\quad
  \vect a=\ddot x\uvec e_x+\ddot y\uvec e_y+\ddot z\uvec e_z.
  \label{r2ch01:eq:009}
\end{equation}
路程 $s$ 是标量，满足 $\dot s=v=\lvert\vect v\rvert$；位移 $\Delta\vect r$ 是矢量，一般有 $\lvert\Delta\vect r\rvert\leq \Delta s$。

直角基矢固定，使“先投影、再求导”与“先求导、再投影”等价。例如
\[
  \uvec e_x\cdot\frac{\dd\vect r}{\dd t}
  =\frac{\dd}{\dd t}(\uvec e_x\cdot\vect r)=\dot x.
\]
这一交换依赖于 $\dot{\uvec e}_x=0$。在极坐标、自然坐标和转动参考系中，基矢会改变，导数的分量便不等于分量的导数，后面出现的牵连项、向心项和科里奥利项都可追溯到这一点。
同样有 $\uvec e_x\cdot\dd\vect v/\dd t=\dd(\uvec e_x\cdot\vect v)/\dd t=a_x$。
常见的另一组基矢记号是 $\uvec i,\uvec j,\uvec k$，分别对应 $\uvec e_x,\uvec e_y,\uvec e_z$，
于是 $\vect r=x\uvec i+y\uvec j+z\uvec k$。
这组记号与哈密顿（Hamilton）的四元数传统有关；
本章只把它们当作三根单位基矢使用，矢量的点积、叉积与四元数乘法不能不加区分地互换。

解析几何中，一条光滑空间曲线可用参数方程
\[
  x=x(\lambda),\qquad y=y(\lambda),\qquad z=z(\lambda)
\]
定义。参数不必是时间；同一几何曲线可由很多不同参数化给出。把 $\lambda$ 选为时间，才得到具体运动。因而“轨迹方程”只说明经过哪些点，“运动方程”还说明何时经过这些点。两个质点可能具有完全相同的轨迹，却以不同快慢、不同方向运动。

匀强重力场中的一般矢量解为
\[
  \vect r(t)=\vect r_0+\vect v_0t+\frac12\vect g t^2,
\]
它是微分方程 $\ddot{\vect r}=\vect g$ 在初值 $\vect r(0)=\vect r_0$、$\dot{\vect r}(0)=\vect v_0$ 下的唯一解。式~\eqref{r2ch01:eq:010} 是把这个矢量式投影到水平、竖直方向后的结果。

分量法不仅是书写方便，还给出了求解策略。若力学规律能分别确定 $a_x,a_y,a_z$，便可对每个标量方程积分，再用同一组初始条件合成。反过来，如果已知矢量运动方程，任意固定方向 $\uvec n$ 上的投影满足
\[
  v_n=\frac{\dd}{\dd t}(\uvec n\cdot\vect r),\qquad
  a_n=\frac{\dd^2}{\dd t^2}(\uvec n\cdot\vect r).
\]
但“固定方向”四字不可省略。若 $\uvec n$ 随时间转动，乘积求导还要出现
$\dot{\uvec n}\cdot\vect r$，这正是活动标架公式的起点。

参数曲线的光滑性要求相邻点随参数连续变化，并在普通点具有非零切向量。若 $\dot{\vect r}=0$，当前参数化可能在该点退化，不能直接把 $\dot{\vect r}/|\dot{\vect r}|$ 当作切向量；应改用弧长、换参数或考察最低阶非零导数。抛体在最高点只有竖直速度分量为零，总速度通常不为零，轨迹依然光滑；竖直上抛最高点则速度完全为零，但几何轨迹仍是一条直线。
在没有自交的局部正则曲线段上，参数区间与曲线段可以一一连续对应，且反向对应也连续，
这种关系称为同胚（homeomorphism）。它比“给定一个参数映射”要求更强；
一条整闭合曲线或含自交点的整条轨迹，不能不加限制地当作与实轴区间同胚。

路程和位移的差异也可从参数曲线看出：
\[
  s(t_2)-s(t_1)=\int_{t_1}^{t_2}|\dot{\vect r}(t)|\,\dd t,\qquad
  \Delta\vect r=\int_{t_1}^{t_2}\dot{\vect r}(t)\,\dd t.
\]
前者把每一小段长度按正值累加，后者按方向矢量相加。只有整段运动方向不改变时，路程才等于位移大小。速度为位置对时间的导数，速率则为路程对时间的导数，二者不可混用。

速度空间中的“位置”是一个速度矢量，所以其坐标量纲为 $LT^{-1}$；速度空间曲线对时间的切向量是加速度，量纲为 $LT^{-2}$。若改用动量 $\vect p=m\vect v$ 作为坐标，就得到动量空间，曲线切向量是合力。这个方法以后还会用于碰撞、天体运动和带电粒子运动。

以斜抛运动为例，取竖直向上为 $z$ 轴，忽略空气阻力，图~\ref{r2ch01:fig:004} 中
\begin{equation}
  x=v_0\cos\phi\,t,\qquad
  z=v_0\sin\phi\,t-\frac12gt^2.
  \label{r2ch01:eq:010}
\end{equation}
消去时间得到轨迹方程
\begin{equation}
  z=x\tan\phi-\frac{g x^2}{2v_0^2\cos^2\phi}.
  \label{r2ch01:eq:011}
\end{equation}
式~\eqref{r2ch01:eq:010} 给出每一时刻的位置，式~\eqref{r2ch01:eq:011} 只给出几何轨迹，二者不能混同。

\figProjectile

\subsection{极坐标与活动标架}
\label{r2ch01:sec:polar}

平面极坐标中
\begin{equation}
  \vect r=r\uvec e_r,\qquad
  \uvec e_r=(\cos\theta,\sin\theta),\qquad
  \uvec e_\theta=(-\sin\theta,\cos\theta).
  \label{r2ch01:eq:012}
\end{equation}
图~\ref{r2ch01:fig:005} 中的基矢会随 $\theta$ 转动，故
\begin{equation}
  \dot{\uvec e}_r=\dot\theta\uvec e_\theta,\qquad
  \dot{\uvec e}_\theta=-\dot\theta\uvec e_r.
  \label{r2ch01:eq:013}
\end{equation}

\figPolarFrame

极坐标由极点、极轴、矢径 $r$ 和极角 $\theta$ 构成。只要轨迹不穿过极点，通常可以把 $r(t),\theta(t)$ 取成光滑函数；$\theta$ 允许超出 $[0,2\pi)$，以保持转动过程连续。位置矢量对两个坐标的偏导分别沿径向和横向：
\[
  \frac{\partial\vect r}{\partial r}=\uvec e_r,\qquad
  \frac{\partial\vect r}{\partial\theta}=r\uvec e_\theta.
\]
所以微小位移
\[
  \dd\vect r=\uvec e_r\,\dd r+r\uvec e_\theta\,\dd\theta.
\]
坐标增量 $\dd\theta$ 本身不是长度，横向实际弧长为 $r\dd\theta$。弧长与坐标增量的比值称为拉梅系数；平面极坐标的两个拉梅系数分别是 $1,r$。柱坐标和球坐标的速度、加速度公式也可由各方向的拉梅系数与活动基矢求导得到。
图~\ref{r2ch01:fig:005}(b) 把这些微分量放大：$PQ$ 沿当前位置的径向，长度为 $\dd R=\dd r$；
$QP'$ 沿当前横向，长度为 $\dd\Theta=r\dd\theta$。两段矢量相加得到斜边 $PP'$，即 $\dd\vect r$。
这里的 $P'$ 是一阶微分构造中的相邻点；有限转角还会产生高阶修正，不能拿放大的三角形反算有限角度。
基矢本身也在转动：若 $\dd\vect\theta=\dd\theta\,\uvec e_z$ 指向纸外，则
\[
 \dd\uvec e_r=\dd\vect\theta\times\uvec e_r=\dd\theta\,\uvec e_\theta,
 \qquad
 \dd\uvec e_\theta=\dd\vect\theta\times\uvec e_\theta=-\dd\theta\,\uvec e_r.
\]
这两个符号可用右手定则逐项判断，与图(a) 中径向和横向的约定一致。

上式可以逐项由全微分验算。把 $\vect r(r,\theta)=r\uvec e_r(\theta)$ 看成两个自变量的矢量函数，有
\[
  \dd\vect r=\frac{\partial\vect r}{\partial r}\dd r
  +\frac{\partial\vect r}{\partial\theta}\dd\theta.
\]
第一个偏导数等于 $\uvec e_r$；第二个偏导数的模等于半径 $r$，方向沿增大 $\theta$ 的切向，所以等于 $r\uvec e_\theta$。因此“坐标变化量”和“真实位移分量”必须通过拉梅系数联系，不能把 $\dot\theta$ 直接当成切向速度。

对活动标架中的矢量 $\vect A=A_r\uvec e_r+A_\theta\uvec e_\theta$，总变化可分为分量在标架内的相对变化
\[
  \dd'\vect A=\dd A_r\uvec e_r+\dd A_\theta\uvec e_\theta
\]
和基矢旋转带来的牵连变化
\[
  \widetilde{\dd}\vect A=A_r\dd\uvec e_r+A_\theta\dd\uvec e_\theta.
\]
由于 $\dd\uvec e_r=\dd\theta\,\uvec e_\theta$、$\dd\uvec e_\theta=-\dd\theta\,\uvec e_r$，两部分相加便得到式~\eqref{r2ch01:eq:014}。这说明导数算符与向活动基矢投影并不对易：
\[
  \uvec e_r\cdot\partial_t\vect A\ne
  \partial_t(\uvec e_r\cdot\vect A).
\]

对任意平面矢量 $\vect A=A_r\uvec e_r+A_\theta\uvec e_\theta$，求导时必须同时对分量和基矢求导：
\begin{equation}
  \left(\frac{\dd\vect A}{\dd t}\right)_r=\dot A_r-\dot\theta A_\theta,\qquad
  \left(\frac{\dd\vect A}{\dd t}\right)_\theta=\dot A_\theta+\dot\theta A_r.
  \label{r2ch01:eq:014}
\end{equation}
于是速度和加速度为
\begin{align}
  \vect v&=\dot r\uvec e_r+r\dot\theta\uvec e_\theta,\label{r2ch01:eq:015}\\
  \vect a&=(\ddot r-r\dot\theta^2)\uvec e_r
  +(r\ddot\theta+2\dot r\dot\theta)\uvec e_\theta.\label{r2ch01:eq:016}
\end{align}
第二式中 $-r\dot\theta^2$ 是向心项，$2\dot r\dot\theta$ 是活动标架转动与径向运动共同产生的项。把它误删，是极坐标运动学最常见的错误之一。

继续对加速度求导，可得加加速度分量
\[
  j_r=\dddot r-3\dot\theta\ddot\theta r-3\dot\theta^2\dot r,
\]
\[
  j_\theta=\dddot\theta r+3\ddot\theta\dot r
  +3\dot\theta\ddot r-\dot\theta^3r.
\]
高阶公式看似复杂，实质仍是反复使用“分量变化加基矢变化”。在实际计算中，先写清基矢导数，通常比死记公式可靠。

斜抛运动也可从发射点用极坐标表示。由式~\eqref{r2ch01:eq:010} 可得
\begin{equation}
  r=\frac{2v_0^2\cos^2\phi}{g}\,
  \frac{\tan\phi-\tan\theta}{\cos\theta},
  \label{r2ch01:eq:017}
\end{equation}
或先由 $t=x/(v_0\cos\phi)$ 求出 $r(t)$ 与 $\theta(t)$。判断质点是否始终远离发射点，应检查 $\dot r=\vect v\cdot\uvec e_r$ 的符号，而不能只看水平分量。

把这个转换实际写开，取 $t>0$，有
\[
 r^2=x^2+z^2=v_0^2t^2-gv_0\sin\phi\,t^3+\frac14g^2t^4,
 \qquad \tan\theta=\frac{z}{x}=\tan\phi-\frac{gt}{2v_0\cos\phi}.
\]
若另选竖直正向 $z$ 为极轴，以 $\theta_z$ 表示从 $z$ 朝 $x$ 量起的角，则
\[
 \theta_z=\frac\pi2-\theta,\qquad
 \tan\theta_z=\frac{x}{z}=\frac{2v_0\cos\phi}{2v_0\sin\phi-gt}.
\]
正切只确定模 $\pi$ 的角，必须结合 $x,z$ 的符号连续选取象限；$z=0$ 时应直接使用位置向量而非除以 $z$。
两套角度只是极轴不同，不能把一个式子的 $\theta_z$ 误代入另一个式子的 $\theta$。

判断远离时，使用乘积求导比直接求根式导数简洁：在 $r>0$ 时，
\[
 r\dot r=\frac{\dd}{\dd t}\left(\frac12r^2\right)
 =\vect r\cdot\vect v
 =t\left(v_0^2-\frac32gv_0\sin\phi\,t+\frac12g^2t^2\right).
\]
对于 $0<\phi<\pi/2$，飞行区间为 $0<t\leq2v_0\sin\phi/g$；括号内二次式的顶点
$t_*=3v_0\sin\phi/(2g)$ 落在这个区间内，最小值为
$v_0^2(1-9\sin^2\phi/8)$。因此始终有 $\dot r>0$ 的条件是
$\sin\phi<2\sqrt2/3$；允许某个瞬间 $\dot r=0$ 而距离不减时取等号。
同一判据也可从该二次式的判别式
$(3gv_0\sin\phi/2)^2-4(g^2/2)v_0^2<0$ 得到。
习题~\ref{r2ch01:pr:001} 将继续用这一关系求完整极坐标轨迹。

\clearpage
\subsection{追逐曲线}
\label{r2ch01:sec:pursuit}

\begin{examplebox}{正方形顶点上的四质点追逐}
四个速率均为 $v$ 的质点最初位于正方形四个顶点，每个质点始终追向顺次相邻的质点。
设初始时顶点到中心的距离为 $r_0$。确定质点的运动方程、轨迹和共同相遇的时刻。
\label{r2ch01:ex:001}
\begin{solution}
对称性保证任意时刻四点仍构成一个转动并缩小的正方形。设顶点到中心的距离为 $r$，相邻方向相差 $\pi/2$，每个速度在径向与切向的投影大小相等。
图~\ref{r2ch01:fig:006}(b) 中 $P$ 追向 $Q$，$\vect v$ 沿 $PQ$；
$\uvec e_r$ 沿 $CP$ 向外，$\uvec e_\theta$ 由它逆时针转过直角。
正方形的几何关系给出 $\vect v$ 与 $\uvec e_r$ 夹角 $3\pi/4$，因此
\begin{equation}
  \dot r=-\frac{v}{\sqrt2},\qquad
  r\dot\theta=\frac{v}{\sqrt2},\qquad
  \frac{\dd r}{\dd\theta}=-r.
  \label{r2ch01:eq:018}
\end{equation}
积分得
\begin{equation}
  r=r_0\e^{-(\theta-\theta_0)},
  \label{r2ch01:eq:019}
\end{equation}
轨迹是对数螺线，如图~\ref{r2ch01:fig:006}。相遇时间为 $\sqrt2 r_0/v$。该推导的关键不是直接猜曲线，而是先利用对称性确认正方形结构始终保持。

也可先分别积分时间方程：
\[
  r(t)=r_0-\frac{vt}{\sqrt2},\qquad
  \theta(t)=\theta_0-\ln\left(1-\frac{vt}{\sqrt2r_0}\right).
\]
消去 $t$ 就回到式~\eqref{r2ch01:eq:019}。直角坐标中只能直接写出
$\dot x^2+\dot y^2=v^2$ 以及位矢、速度夹角为 $3\pi/4$，联立后虽然也能求解，却掩盖了径向、切向速度恒定比例的简单结构。选择坐标系的意义就在于把已知条件放到最自然的分量上。
例如，在图示逆时针追逐的约定下，还可写成
\[
 \frac{\vect r\cdot\vect v}{|\vect r\times\vect v|}
 =\frac{x\dot x+y\dot y}{x\dot y-y\dot x}=-1.
\]
这对应 $\cot(3\pi/4)=-1$；一般角度下点积与叉积模之比是余切，不是正切。
也可以尝试直接联立这两条直角坐标约束求解，并与极坐标结果互相检查，体会选系如何简化方程。

\figPursuitFour
\end{solution}
\end{examplebox}

\subsection{自然坐标与曲率}
\label{r2ch01:sec:frenet}

自然坐标的构造属于古典微分几何。牛顿等人发展出的微积分、笛卡尔等人建立的解析几何，
使研究曲线和曲面的局部性质可以转化为函数、导数与坐标之间的关系。

自然坐标的出发点是：不再用事先固定的 $x,y,z$ 方向分解运动，而让基矢由轨迹在当前点的几何性质决定。取曲线上一点 $P$，连接邻近点 $Q$ 得割线 $PQ$；当 $Q$ 从曲线两侧趋近 $P$ 时，割线的极限方向是切线方向。再取两个邻近点 $Q,R$，三点 $P,Q,R$ 确定一个圆；当二者同时趋近 $P$ 时，若极限存在，所得圆就是密切圆。它不仅经过 $P$ 且切线相同，还在二阶弯曲程度上与原曲线一致。

密切圆的圆心 $C$ 是曲率中心，$\overrightarrow{PC}$ 给出主法向；半径 $\rho=PC$ 是曲率半径。主法向不一定指向某个全局“内侧”，而是指向切向量变化的方向。切向量与主法向张成密切平面，副法向 $\uvec\beta=\uvec\tau\times\uvec n$ 垂直于密切平面。对平面曲线，密切平面固定；对空间曲线，它可沿轨迹扭转。

用弧长 $s$ 作参数具有特殊优势。普通参数 $\lambda$ 下，
$\dd\vect r/\dd\lambda$ 的模依赖参数走得多快；弧长参数下
$|\dd\vect r/\dd s|=1$，导数只表示方向。若曲线有尖点或参数化在某点退化，自然标架可能不存在或需要分段定义；曲率公式的分母为零时也必须回到几何极限判断，而不能盲目代数化简。

用弧长 $s$ 参数化光滑曲线，定义单位切向量、主法向量和副法向量
\begin{equation}
  \uvec\tau=\frac{\dd\vect r}{\dd s},\qquad
  \frac{\dd\uvec\tau}{\dd s}=\kappa\uvec n,\qquad
  \uvec\beta=\uvec\tau\times\uvec n.
  \label{r2ch01:eq:020}
\end{equation}
其中 $\kappa=1/\rho$ 为曲率，$\rho$ 为曲率半径。三维曲线的弗雷内公式为
\begin{equation}
  \frac{\dd\uvec\tau}{\dd s}=\kappa\uvec n,\qquad
  \frac{\dd\uvec n}{\dd s}=-\kappa\uvec\tau+\tau_g\uvec\beta,\qquad
  \frac{\dd\uvec\beta}{\dd s}=-\tau_g\uvec n,
  \label{r2ch01:eq:021}
\end{equation}
这里 $\tau_g$ 是曲线的挠率。图~\ref{r2ch01:fig:007} 中，$\uvec\tau$、$\uvec n$ 张成密切平面，密切圆圆心沿 $\uvec n$ 方向，半径就是 $\rho$。

为避免只记结果，可以把三个导数先写成最一般形式
\[
\begin{aligned}
  \uvec\tau'&=a_1\uvec\tau+a_2\uvec n+a_3\uvec\beta,\\
  \uvec n'&=b_1\uvec\tau+b_2\uvec n+b_3\uvec\beta,\\
  \uvec\beta'&=c_1\uvec\tau+c_2\uvec n+c_3\uvec\beta,
\end{aligned}
\qquad({}'=\dd/\dd s).
\]
对 $\uvec\tau^2=\uvec n^2=\uvec\beta^2=1$ 分别求导得 $a_1=b_2=c_3=0$；对三组正交关系求导得
$a_2=-b_1$、$a_3=-c_1$、$b_3=-c_2$。于是九个系数只剩三个独立量。再利用主法向就按 $\uvec\tau'$ 方向定义，可知 $a_3=c_1=0$；剩下的 $a_2=\kappa$、$b_3=\tau_g$ 就给出式~\eqref{r2ch01:eq:021}。这个推导同时说明弗雷内标架的旋转率矩阵必然反对称。

\figFrenet

图~\ref{r2ch01:fig:007} 将点 $P$ 附近的密切平面画成纸面，曲率中心 $C$ 位于曲线的凹侧。
半径 $CP$ 垂直于 $P$ 处切线，所以 $\uvec n$ 指向 $C$，而 $\uvec\tau$ 与 $\uvec n$ 的右手叉积
$\uvec\beta$ 指向纸面外，用 $\odot$ 表示。$Q,R$ 是曲线上的邻近点；它们与 $P$ 确定的圆在
$Q,R\to P$ 时才趋于图中的密切圆，不要求有限距离处的 $Q,R$ 已经落在该圆上。

自然坐标不是由预先选定的背景坐标轴决定，而是由轨迹自身的局部几何决定。规定曲线上的起点后，每一点可用从起点量起的弧长 $s$ 标记，因而曲线本身写成 $\vect r=\vect r(s)$。由于 $\lvert\dd\vect r\rvert=\dd s$，切向量 $\dd\vect r/\dd s$ 自动是单位矢量。取曲线上邻近点 $Q,R$ 与 $P$ 构成的圆，当 $Q,R\to P$ 时，其极限是密切圆；圆心指向定义主法向，圆所在平面是密切平面，副法向垂直于该平面。

弗雷内公式可由正交归一性推出，而不必当作三个孤立结论记忆。设
\[
  \frac{\dd\uvec\tau}{\dd s}=a_1\uvec\tau+a_2\uvec n+a_3\uvec\beta,
\]
并对 $\uvec\tau\cdot\uvec\tau=1$ 求导，可知 $a_1=0$。主法向定义为 $\dd\uvec\tau/\dd s$ 的方向，所以 $a_3=0$，剩下的正系数就是曲率 $\kappa$。再分别对
$\uvec n\cdot\uvec n=1$、$\uvec\beta\cdot\uvec\beta=1$ 以及三组互相正交条件求导，各导数矩阵必须是反对称矩阵，由此得到式~\eqref{r2ch01:eq:021}。曲率表示沿曲线前进单位弧长时切线转过的角度：
\[
  \kappa=\left|\frac{\dd\uvec\tau}{\dd s}\right|
  =\left|\frac{\dd\vartheta}{\dd s}\right|\quad\text{（平面曲线）}.
\]
挠率 $\tau_g$ 描述密切平面绕切线方向转动的快慢；平面曲线的挠率恒为零。
这里 $\vartheta$ 是平面切线相对于固定方向的带符号角度，取绝对值才能保证曲率非负。
在三维情形，应使用 $|\dd\uvec\tau/\dd s|$ 的定义，不能用一个平面方位角代替空间切向变化。
弗雷内标架要求当前位置速率非零、曲率非零；在直线段或拐点处，主法向与副法向需另作连续选取，不能直接除以零曲率。

还可以从一般反对称转率矩阵理解自然标架。三个单位基矢的导数在自身方向投影都为零；任意两基矢保持垂直，要求一个基矢导数在另一个方向的投影与反向投影正负相消。最一般形式本来含有曲率、挠率以及切向量导数的副法向分量；主法向正是按切向量导数方向定义，最后一个分量因此被定义为零。弗雷内公式中的两个几何标量 $\kappa,\tau_g$ 已经完整决定标架如何沿曲线转动。

曲率 $\kappa$ 的量纲为 $L^{-1}$，挠率同样为 $L^{-1}$。直线的曲率为零、曲率半径无穷大；半径 $R$ 的圆有常曲率 $1/R$、挠率为零；圆柱螺线既有曲率又有挠率。曲率只反映弯得多急，挠率只反映弯曲平面如何扭转，它们都与质点沿曲线的实际速率无关。

在 $s=0$ 的邻域，曲线可作局部展开
\begin{equation}
  \vect r(s)=\vect r_0+s\uvec\tau_0
  +\frac12\kappa_0s^2\uvec n_0
  +\frac16s^3\bigl(-\kappa_0^2\uvec\tau_0+\kappa_0'\uvec n_0+\kappa_0\tau_{g0}\uvec\beta_0\bigr)
  +\order{s^4}.
  \label{r2ch01:eq:022}
\end{equation}
该展开给出自然坐标系中的局部形状。若只保留最低非零阶，切向、法向、副法向坐标分别为
\[
  T=s+\order{s^3},\qquad
  N=\frac12\kappa s^2+\order{s^3},\qquad
  B=\frac16\kappa\tau_gs^3+\order{s^4}.
\]
因此曲线在足够小的邻域内首先像一条切线，二阶才显现平面内弯曲，三阶才显现离开密切平面的扭转。密切圆正是保留到二阶时与原曲线贴合程度最高的圆。

这个阶次结构也解释了为什么普通动力学最常用速度和加速度，却较少需要挠率：速度是一阶时间导数，只看切线；加速度是二阶，只看切向速率变化和曲率法向；要让副法向出现，至少要到三阶加加速度。式~\eqref{r2ch01:eq:025} 中的
$\kappa\tau_gv^3\uvec\beta$ 正是空间轨迹扭转的第一个运动学信号。

若再把弧长作为时间函数展开，
\[
  \Delta s=\dot s\,\Delta t+\frac12\ddot s\,\Delta t^2
  +\frac16\dddot s\,\Delta t^3+\cdots,
\]
并代入曲线的弧长展开，就得到位置的时间泰勒式
\[
  \Delta\vect r=\vect v\Delta t+\frac12\vect a\Delta t^2
  +\frac16\vect j\Delta t^3+\cdots.
\]
式~\eqref{r2ch01:eq:023}—\eqref{r2ch01:eq:025} 正是逐阶比较所得。速度只有切向分量；加速度首次出现法向分量；副法向分量直到加加速度才通过挠率出现。
由 $\dot s=v$ 得
\begin{align}
  \vect v&=v\uvec\tau,\label{r2ch01:eq:023}\\
  \vect a&=\dot v\uvec\tau+\kappa v^2\uvec n,\label{r2ch01:eq:024}\\
  \vect j&=(\ddot v-\kappa^2v^3)\uvec\tau
  +(3\kappa v\dot v+\dot\kappa v^2)\uvec n
  +\kappa\tau_gv^3\uvec\beta.\label{r2ch01:eq:025}
\end{align}
因此不需要先求轨迹方程也可由速度、加速度求曲率半径：
\begin{equation}
  \kappa=\frac{\lvert\vect v\times\vect a\rvert}{v^3},
  \qquad
  \rho=\frac{v^3}{\lvert\vect v\times\vect a\rvert}.
  \label{r2ch01:eq:026}
\end{equation}
平面曲线 $y=y(x)$ 的曲率为
\begin{equation}
  \kappa=\frac{|y''|}{(1+y'^2)^{3/2}}.
  \label{r2ch01:eq:027}
\end{equation}
极坐标曲线 $r=r(\theta)$ 则有
\begin{equation}
  \rho=\frac{\bigl(r^2+r'^2\bigr)^{3/2}}
  {\left|r^2+2r'^2-r r''\right|}.
  \label{r2ch01:eq:028}
\end{equation}
例如对数螺线 $r=r_0\e^{k\theta}$ 满足 $\rho=\sqrt{1+k^2}\,r$。

以 $y=f(x)$ 验证式~\eqref{r2ch01:eq:027}：为了研究几何曲率，可以让一个辅助质点
以任意常量横向速度 $u\ne0$ 沿曲线运动，即 $x=ut$、$y=f(ut)$。于是
\[
 \vect v=(u,uf'),\qquad \vect a=(0,u^2f''),\qquad
 |\vect v\times\vect a|=|u|^3|f''|.
\]
将速率 $v=|u|\sqrt{1+f'^2}$ 代入式~\eqref{r2ch01:eq:026}，
$|u|^3$ 消去，得到 $\rho=(1+f'^2)^{3/2}/|f''|$。选择不同的 $u$ 不改变曲率，
也不必把量纲不同的坐标 $x$ 和时间 $t$ 直接当成相等。

这个结论也可用运动学直接验证。令 $\theta=\omega t$，则
\[
  \vect v=(k\omega r,\omega r),\qquad
  \vect a=((k^2-1)\omega^2r,2k\omega^2r)
\]
是极坐标活动标架中的分量。代入式~\eqref{r2ch01:eq:026} 后，$\omega$ 完全消去，说明曲率是轨迹的几何性质，与人为选择质点沿曲线运动的快慢无关。
具体地，$v=|\omega|r\sqrt{1+k^2}$，并且
\[
 |\vect v\times\vect a|
 =|v_ra_\theta-v_\theta a_r|
 =|\omega|^3r^2(1+k^2),
\]
所以 $\rho=v^3/|\vect v\times\vect a|=r\sqrt{1+k^2}$。这里假设 $r>0$、$\omega\ne0$；
$k=0$ 时轨迹退化为圆，公式给出 $\rho=r$，与圆的几何性质一致。

\section{刚体运动与参考系变换}
\label{r2ch01:sec:rigid}

\subsection{固连标架与角速度}

刚体中任意两点的距离保持不变。取刚体上一点 $C$ 为基点，并建立与刚体固连的正交标架 $(\uvec e_1,\uvec e_2,\uvec e_3)$，如图~\ref{r2ch01:fig:008}。刚体的瞬时运动可唯一分解为基点平动和绕基点转动。角速度 $\vect\omega$ 由
\begin{equation}
  \left(\frac{\dd\vect A}{\dd t}\right)_{S}
  =\left(\frac{\dd\vect A}{\dd t}\right)_{S'}
  +\vect\omega\times\vect A
  \label{r2ch01:eq:029}
\end{equation}
定义；它适用于任意矢量 $\vect A$。特别地，固连基矢在动系中的分量不变，故 $\dot{\uvec e}_i=\vect\omega\times\uvec e_i$。

刚体是一种理想化：真实物体受力会有形变，只要形变在研究精度内可忽略，就可以把各点间距离视为常数。这个条件比“物体形状看起来不变”更严格，也比“选一个点不动”更一般。刚体可以整体平动、绕固定轴转动，也可以在三维空间中同时平移和改变转轴方向。

与刚体固连的标架由一个原点和三根彼此垂直、长度为一的基矢组成。原点运动由 $\vect r_C(t)$ 描述，标架姿态原则上要用三个独立转角描述；角速度 $\vect\omega(t)$ 则给出姿态在当前瞬间的变化率。只知道 $\vect\omega$ 的一个瞬时值不能恢复完整历史姿态，但足以计算该瞬间所有固连点的速度；再知道角加速度，就可计算加速度。

角速度矢量的存在可以从固连标架的正交归一性严格推出。把三个基矢的导数写成
\[
  \dot{\uvec e}_i=\sum_j\Omega_{ji}\uvec e_j.
\]
对 $\uvec e_i\cdot\uvec e_j=\delta_{ij}$ 求导，得到
$\Omega_{ij}+\Omega_{ji}=0$，所以系数矩阵是三阶反对称矩阵。任意三阶反对称矩阵都与唯一矢量 $\vect\omega$ 对应，使它作用在任意矢量上等同于叉乘 $\vect\omega\times$。因此
\[
  \dot{\uvec e}_1=\vect\omega\times\uvec e_1,\quad
  \dot{\uvec e}_2=\vect\omega\times\uvec e_2,\quad
  \dot{\uvec e}_3=\vect\omega\times\uvec e_3.
\]
这不是只适用于某一根杆的局部公式，而是整个刚体共同的瞬时角速度。

若把反对称矩阵写开，三个独立分量可组合成
\[
  \vect\omega=\omega_1\uvec e_1+\omega_2\uvec e_2+\omega_3\uvec e_3.
\]
具体地，
\[
  \begin{pmatrix}
  \dot{\uvec e}_1\\ \dot{\uvec e}_2\\ \dot{\uvec e}_3
  \end{pmatrix}
  =
  \begin{pmatrix}
  0&\omega_3&-\omega_2\\
  -\omega_3&0&\omega_1\\
  \omega_2&-\omega_1&0
  \end{pmatrix}
  \begin{pmatrix}
  \uvec e_1\\ \uvec e_2\\ \uvec e_3
  \end{pmatrix}.
\]
这里一列中排列的是三个\emph{基矢本身}，不是同一矢量的三个分量。若改为对分量列 $(a_1,a_2,a_3)^{\mathsf T}$ 计算 $\vect\omega\times\vect a$，矩阵应为
\[
  [\vect\omega\times]=
  \begin{pmatrix}
  0&-\omega_3&\omega_2\\
  \omega_3&0&-\omega_1\\
  -\omega_2&\omega_1&0
  \end{pmatrix},
\]
恰好是上一个矩阵的转置；两种写法的行列含义不同，不能直接互换。基矢导数矩阵的对角元为零表示单位基矢的导数不能沿自身方向；非对角元成对反号，则保证任意两根基矢始终正交。
例如
\[
  \dot{\uvec e}_1=\omega_3\uvec e_2-\omega_2\uvec e_3,
\]
另外两式循环得到。基矢对时间的一阶导数只有三个独立自由度，这与三维转动恰有三个瞬时自由度一致。$\vect\omega$ 不是某个点独有的量，也不是把三个欧拉角分别求导后简单拼接；它是所有固连矢量共同满足的几何转率。

\figRigidFrame

若某点相对固连标架的三个坐标始终不变，就说该点固连于标架。若物体所有点都能固连于同一个正交标架，任意两点距离便保持不变，这就是刚体模型。反过来，任取刚体上三个不共线的点，以其中一点为原点，以指向第二点的单位矢量为第一基矢，再把指向第三点的矢量去掉第一基矢方向的投影并归一化，得到第二基矢；二者的右手叉积就是第三基矢。这样便构造出固连正交标架。因此“刚体”和“与刚体固连的参考系”在运动学上可以相互定义。

取两个固连点 $C,P$，相对矢量 $\vect\rho=\overrightarrow{CP}$ 在动系中的三个分量都是常数，所以
$\dot{\vect\rho}=\vect\omega\times\vect\rho$。加上基点位置
$\vect r_P=\vect r_C+\vect\rho$，便得到速度基点公式；再求导得到加速度基点公式。整个推导没有假设基点是质心、转轴点或静止点，因此基点可以按解题方便任意选择。

任意固连矢量
\[
  \vect A=A_1\uvec e_1+A_2\uvec e_2+A_3\uvec e_3
\]
在动系中的分量 $A_i$ 为常数，在定系中仍会因基矢旋转而变化：
\[
  \dot{\vect A}=\sum_iA_i\dot{\uvec e}_i
  =\vect\omega\times\vect A.
\]
这条公式与矢量的量纲无关：位移、速度、力、角动量等只要相对刚体的大小和方向均固定，即固连分量恒定，都服从同一转动导数。
若只是方向固定而大小 $A(t)$ 改变，则还须保留 $\dot A\,\uvec A$ 这一相对导数项，不能只写叉积。
$\vect\omega$ 的方向给出瞬时转轴方向，大小给出单位时间转角。

设刚体上一点 $P$ 相对 $C$ 的位置为 $\vect\rho$，则
\begin{align}
  \vect v_P&=\vect v_C+\vect\omega\times\vect\rho,
  \label{r2ch01:eq:030}\\
  \vect a_P&=\vect a_C+\vect\beta\times\vect\rho
  +\vect\omega\times(\vect\omega\times\vect\rho),
  \qquad \vect\beta=\dot{\vect\omega}.
  \label{r2ch01:eq:031}
\end{align}
式~\eqref{r2ch01:eq:030} 中相对速度垂直于 $\vect\rho$，所以刚体任意两点 $P,Q$ 必有
\begin{equation}
  (\vect v_P-\vect v_Q)\cdot(\vect r_P-\vect r_Q)=0.
  \label{r2ch01:eq:032}
\end{equation}
这条关系常用于由若干点的速度反求角速度。

加速度式中的三项具有不同几何意义：$\vect a_C$ 是基点平动，$\vect\beta\times\vect\rho$ 是由角速度大小或方向改变产生的切向项，$\vect\omega\times(\vect\omega\times\vect\rho)$ 指向瞬时转轴，是向轴加速度。选择不同基点会改变三项各自的数值，但总和 $\vect a_P$ 不变。

这两条基点公式可以从位置关系一步不略地推出。先写
\[
  \vect r_P=\vect r_C+\vect\rho,
  \qquad \left(\frac{\dd\vect\rho}{\dd t}\right)_{\text{刚体系}}=0.
\]
固连矢量在定系中的导数为 $\dot{\vect\rho}=\vect\omega\times\vect\rho$，所以第一次求导得
$\vect v_P=\vect v_C+\vect\omega\times\vect\rho$。第二次求导时必须对叉积的两个因子都求导：
\[
\begin{aligned}
  \vect a_P
  &=\vect a_C+\dot{\vect\omega}\times\vect\rho
    +\vect\omega\times\dot{\vect\rho}\\
  &=\vect a_C+\vect\beta\times\vect\rho
    +\vect\omega\times(\vect\omega\times\vect\rho).
\end{aligned}
\]
向轴项中出现两次 $\vect\omega$，正是因为相对位矢虽在刚体系内不变，却在定系中以同一角速度改变方向。若把 $\dot{\vect\rho}$ 错设为零，就会遗漏整个向轴加速度。

\subsection{平动、转动参考系中的速度和加速度}
\label{r2ch01:sec:transform}

定系 $S$ 的原点为 $O$，动系 $S'$ 的原点为 $C$，同一点 $P$ 的位置满足图~\ref{r2ch01:fig:009} 所示矢量三角形
\begin{equation}
  \vect r=\vect r_C+\vect r'.
  \label{r2ch01:eq:033}
\end{equation}

\figFrameTransform

同一个相对位矢 $\vect r'$ 可以有三种不同意义的时间变化。绝对导数 $D$ 固定在定系中比较矢量；相对导数 $D'$ 固定在动系中比较分量；随体导数 $\widetilde D$ 只保留动系基矢转动造成的变化。若
\[
  \vect A=A_1\uvec e_1+A_2\uvec e_2+A_3\uvec e_3,
\]
则
\[
  D'\vect A=\dot A_1\uvec e_1+\dot A_2\uvec e_2+\dot A_3\uvec e_3,
  \qquad
  \widetilde D\vect A=\vect\omega\times\vect A,
\]
并且
\[
  D=\widetilde D+D'.
\]
这就是转动求导公式~\eqref{r2ch01:eq:029} 的算符写法。相同符号“对时间求导”若没有说明固定哪个参考系，会造成概念混乱。

三种导数可用三个典型矢量辨认。若 $\vect A$ 固连动系，则 $D'\vect A=0$，但
$D\vect A=\vect\omega\times\vect A$；若 $\vect A$ 在定系中大小、方向不变，则 $D\vect A=0$，它在动系中的分量反而以相反角速度变化；一般矢量同时具有分量变化和基矢变化，两者相加才是绝对变化。极坐标中的“分量的导数不等于导数的分量”就是这条规则在二维旋转基矢上的具体表现。

对式~\eqref{r2ch01:eq:033} 用式~\eqref{r2ch01:eq:029} 求导，得到速度合成公式
\begin{equation}
  \vect v=\vect v_C+\vect\omega\times\vect r'+\vect v',
  \qquad
  \vect v'=\left(\frac{\dd\vect r'}{\dd t}\right)_{S'}.
  \label{r2ch01:eq:034}
\end{equation}
右边三项依次为基点速度、牵连速度和相对速度。再求一次导数：
\begin{align}
  \vect a={}&\vect a_C+\vect\beta\times\vect r'
  +\vect\omega\times(\vect\omega\times\vect r')
  +2\vect\omega\times\vect v'+\vect a',
  \label{r2ch01:eq:035}
\end{align}
其中 $\vect a'=(\dd\vect v'/\dd t)_{S'}$。五项依次是基点加速度、欧拉加速度、向心加速度、科里奥利加速度和相对加速度。若动系仅作平动，$\vect\omega=0$，便退化为
\begin{equation}
  \vect v=\vect v_C+\vect v',
  \qquad \vect a=\vect a_C+\vect a'.
  \label{r2ch01:eq:036}
\end{equation}

为看清两个科里奥利贡献，可把第二次求导完整展开：
\begin{align*}
  D\vect v
  &=
  D\vect v_C+D\vect\omega\times\vect r'
  +\vect\omega\times D\vect r'+D\vect v'\\
  &=
  \vect a_C+\vect\beta\times\vect r'
  +\vect\omega\times(\vect\omega\times\vect r'+\vect v')\\
  &\quad+(\vect\omega\times\vect v'+\vect a')\\
  &=\vect a_C+\vect\beta\times\vect r'
  +\vect\omega\times(\vect\omega\times\vect r')
  +2\vect\omega\times\vect v'+\vect a'.
\end{align*}
其中一项来自 $D\vect r'$ 中的相对速度，另一项来自 $D\vect v'$ 时活动基矢的旋转。整理后正好是式~\eqref{r2ch01:eq:035}。若只把速度公式中每个符号“机械加点”，而忘记 $\vect r'$、$\vect v'$ 都在转动基底中表示，就会少掉一份甚至两份科里奥利项。

牵连点不是一个始终和动点相同的物质点。它只在研究瞬间与动点占据同一空间位置，却固连于动系；下一瞬间二者一般分离。正因为如此，牵连点的速度、加速度可以由动系刚体运动直接求出，而动点相对它的变化就是相对量。机构题、接触题中若不区分这两个点，常会把几何交点速度错当成某个材料点速度。

把
\[
  \vect v_{\mathrm{tr}}=\vect v_C+\vect\omega\times\vect r',
\]
\[
  \vect a_{\mathrm{tr}}=\vect a_C+\vect\beta\times\vect r'
  +\vect\omega\times(\vect\omega\times\vect r')
\]
称为牵连速度和牵连加速度，则速度合成是“绝对等于牵连加相对”，加速度合成还要加
$2\vect\omega\times\vect v'$。科里奥利项有两份来源：相对位矢改变会改变牵连速度一次，相对速度矢量本身又被动系旋转一次。它非零还要求相对速度具有垂直于转轴的分量；若 $\vect\omega=0$、$\vect v'=0$ 或二者平行，叉积均为零。这个项不显含相对位置 $\vect r'$。

使用动点—动系法时，可按以下顺序组织：
\begin{enumerate}
  \item 选定定系、动系和研究的动点；再指定“此刻与动点重合、但始终固连于动系”的牵连点。
  \item 在定系中写动点的绝对速度、加速度；在动系中写相对量；把牵连点当作动系刚体上的点求牵连量。
  \item 用式~\eqref{r2ch01:eq:034}、\eqref{r2ch01:eq:035} 合成，并单独检查科里奥利项的方向。
  \item 最后核验选系是否真的简化了约束。计算本身不要求动系坐标轴一定取直角、极坐标或自然坐标；这些只是分量表达方式。
\end{enumerate}

若 $\vect\omega=\vect\beta=\vect0$，动系是平动系，所有固连直线始终平行，牵连量等于动系任一固定点的平动量。此时没有向心、欧拉和科里奥利项，式~\eqref{r2ch01:eq:036} 即完整结果。

两个固连于不同刚体的参考系也可用同样方法比较。设 $S_2$ 相对 $S_1$ 的原点速度、角速度分别为 $\vect V_{21},\vect\Omega_{21}$，则对同一刚体的角速度
\begin{equation}
  \vect\omega_{B/1}=\vect\omega_{B/2}+\vect\Omega_{21}.
  \label{r2ch01:eq:037}
\end{equation}
角加速度求导时还要考虑动系转动：
\begin{equation}
  \vect\beta_{B/1}=\vect\beta_{B/2}+\vect\beta_{21}
  +\vect\Omega_{21}\times\vect\omega_{B/2}.
  \label{r2ch01:eq:038}
\end{equation}
为证明角速度合成，取任意固连于刚体 $B$ 的矢量 $\vect b$。分别应用转动求导公式，
\[
 \left(\frac{\dd\vect b}{\dd t}\right)_1
 =\left(\frac{\dd\vect b}{\dd t}\right)_2+\vect\Omega_{21}\times\vect b
 =(\vect\omega_{B/2}+\vect\Omega_{21})\times\vect b.
\]
这一式对刚体内任意方向的固连矢量都成立，故其共同角速度为式~\eqref{r2ch01:eq:037}。
再在 $S_1$ 中对它求导，并对 $\vect\omega_{B/2}$ 应用同一转动求导公式：
\[
 \left(\frac{\dd\vect\omega_{B/2}}{\dd t}\right)_1
 =\underbrace{\left(\frac{\dd\vect\omega_{B/2}}{\dd t}\right)_2}_{\vect\beta_{B/2}}
  +\vect\Omega_{21}\times\vect\omega_{B/2}.
\]
加上 $\dd\vect\Omega_{21}/\dd t=\vect\beta_{21}$，即得式~\eqref{r2ch01:eq:038}。
最后一项正是相对角速度矢量被动系转动带来的方向变化率。

若还要变换刚体基点的线运动，设其在 $S_2$ 中的位置、速度、加速度为
$\vect R',\vect V',\vect A'$，则
\[
  \vect V=\vect v_C+\vect\Omega_{21}\times\vect R'+\vect V',
\]
\[
  \vect A=\vect a_C+\vect\beta_{21}\times\vect R'
  +\vect\Omega_{21}\times(\vect\Omega_{21}\times\vect R')
  +2\vect\Omega_{21}\times\vect V'+\vect A'.
\]
线运动公式与点变换完全相同；角运动则由式~\eqref{r2ch01:eq:037}、\eqref{r2ch01:eq:038} 给出。这样，刚体相对刚体的复合运动就被平移和转动两组公式完整描述。

\section{接触、纯滚动与几何约束}
\label{r2ch01:sec:contact}

\subsection{接触点的三个身份}

两个物体接触时，必须区分物体1上此刻位于接触处的物质点 $P_1$、物体2上的物质点 $P_2$，以及接触位置本身这个几何点 $P_3$。图~\ref{r2ch01:fig:010} 中，纯滚圆盘的 $P_1$ 瞬时速度为零，地面的 $P_2$ 始终静止，但几何切点 $P_3$ 会沿地面移动。把 $P_1$ 与 $P_3$ 混为同一点，就会错误地推出“接触点不移动”。

可以用“身份是否随时间更换”检查这三个点。$P_1$ 是圆盘上某个已做标记的材料点，运动后仍是同一块材料，其轨迹为旋轮线；$P_2$ 也是地面上一个固定材料点；$P_3$ 却由“每一时刻都取圆盘最低点”这条几何规则定义，因而每刻都换成圆盘上的新材料点。三者在某一瞬间坐标相同，不代表它们的速度或加速度必然相同。

\figRollingPoints

一般的“相交系”中，几何交点相对每个几何对象都可能移动。分别以两个对象为动系，对同一个几何交点写“牵连加相对”，所得绝对速度和加速度必须相同，由此产生约束方程。相切是相交的特殊情形：交点的相对轨迹分别沿两物体轮廓，瞬时相对速度沿公共切线；两物体保持接触还要求法向没有穿透。

以圆盘在地面滚动为例，盘上材料点 $P_1$ 的轨迹是旋轮线：它只在接地瞬间速度为零，随后立即离开地面；地面材料点 $P_2$ 在地面参考系中始终静止；几何最低点 $P_3$ 则沿地面以圆心平移速度移动。若在圆筒内放一只始终跑到最低位置的小动物，它近似代表几何点 $P_3$ 的运动，而不是同一个圆筒材料点。三个点同位却不同身份，是接触题最需要先澄清的概念。
图~\ref{r2ch01:fig:010}(b) 中，半径为 $R$ 的圆盘向右纯滚过顺时针转角 $\chi$。
圆心从 $(0,R)$ 移到 $(R\chi,R)$；初始接地点对应的材料点 $P_1$ 相对圆心的位置
从 $(0,-R)$ 旋转为 $(-R\sin\chi,-R\cos\chi)$。相加得到旋轮线参数式
\[
 x_{P_1}=R(\chi-\sin\chi),\qquad y_{P_1}=R(1-\cos\chi).
\]
它在 $\chi=\pi$ 时达到 $(\pi R,2R)$，在 $\chi=2\pi$ 时重新接地于 $(2\pi R,0)$；
同一时刻几何切点却一直是 $P_3=(R\chi,0)$。若 $\dot\chi$ 恒定，
\[
 \vect v_{P_1}=R\dot\chi(1-\cos\chi,\sin\chi),\qquad
 \vect v_{P_3}=(R\dot\chi,0).
\]
前者在每次接地时为零，后者并不为零，这给出了图(a) 中身份区别的定量检验。

几何交点的相对速度沿轮廓切线，是因为它在每个物体自身坐标中受到“必须留在轮廓上”的约束。对曲线 $F(\vect r)=0$，相对运动满足
$\nabla F\cdot\vect v'=0$，即相对速度垂直于法向梯度。二次求导后出现曲率引起的法向相对加速度，这就是仅用一次速度条件不足以直接写出全部加速度关系的原因。

若分别把指向两个物体内部的法向取为各自正方向，则两个法向彼此反向，不分离条件写成 $v_{1n}=-v_{2n}$；若统一使用同一个公法线 $\uvec n$，则写成
$(\vect v_1-\vect v_2)\cdot\uvec n=0$。这两种形式物理意义相同，不能把符号差异误认为不同结论。加速度层面还要计入几何交点沿弯曲轮廓运动的法向加速度，以及在转动动系中出现的科里奥利项。

不分离条件要求两物体在公法线方向的速度相同；无滑动还要求公共切线方向速度相同，即
\begin{equation}
  (\vect v_1-\vect v_2)\cdot\uvec n=0,
  \qquad
  (\vect v_1-\vect v_2)\cdot\uvec\tau=0.
  \label{r2ch01:eq:039}
\end{equation}
下面讨论平面内两个光滑凸轮廓的外接触。取公共单位切向 $\uvec\tau$，使其在纸面内逆时针转
$90^\circ$ 后成为指向物体1内部的单位法向 $\uvec n_1$；物体2的内法向为 $\uvec n_2=-\uvec n_1$。
两个轮廓的曲率大小记为 $\kappa_1,\kappa_2\geq0$，且 $\kappa_1+\kappa_2>0$。
角速度 $\omega_1,\omega_2$ 均以逆时针为正。纯滚时，接触处两个材料点有同一个瞬时速度
$\vect v_m$；几何切点 $P_3$ 相对它们的切向速度记为 $u\uvec\tau$。于是
\begin{equation}
  u=\frac{\omega_2-\omega_1}{\kappa_1+\kappa_2},\qquad
  \vect v_{P_3}=\vect v_m+u\uvec\tau.
  \label{r2ch01:eq:040}
\end{equation}
注意分式给出的是相对切向速度，不是任意参考系中的 $\vect v_{P_3}$。
若共同给两物体叠加一个平移速度，$\omega_1,\omega_2$ 与曲率都不变，
而 $\vect v_m,\vect v_{P_3}$ 同时增加该平移速度，式~\eqref{r2ch01:eq:040} 仍成立。
若为内接触或采用带符号曲率，应按法向约定重写曲率差，不能直接套用外接触的曲率和。

纯滚的几何定义是：接触过程中，两物体轮廓滚过的弧长相等，没有相对滑动。
用同一空间切向规定带符号弧长，令 $u=\dot s$，则
\[
  \vect v'_{P_3/1}=\dot s\,\uvec\tau,\qquad
  \vect v'_{P_3/2}=\dot s\,\uvec\tau.
\]
对几何切点分别写速度合成并相减，就得到式~\eqref{r2ch01:eq:039} 的切向部分。若切向分量不相等，接触处必有滑动，只能称为滚滑而非纯滚。

“接触点瞬时静止”只在一个物体相对地面纯滚且地面静止时成立。两个都在运动的齿轮、圆柱或曲面纯滚时，两个接触材料点速度彼此相等，却不一定为零；几何切点速度还可能与它们不同。判断纯滚应比较两个材料点在公共切线方向的速度，而不是寻找一个绝对静止点。

式~\eqref{r2ch01:eq:040} 可由公切线的转率直接推出。设公切线相对固定方向的角度为 $\psi$。
几何切点沿轮廓移动会改变轮廓切向，刚体本身的转动也会改变切向。分别从两物体看，
\[
 \dot\psi=\omega_1+\kappa_1u
          =\omega_2-\kappa_2u.
\]
两式相减，得到 $(\kappa_1+\kappa_2)u=\omega_2-\omega_1$；
若只取速率，则 $(\kappa_1+\kappa_2)|u|=|\omega_2-\omega_1|$。

求加速度时，不能把 $\vect v_1(P_3,t)=\vect v_2(P_3,t)$ 的时间导数直接写成
$\vect a_1=\vect a_2$，因为跟踪的是不断换材料的几何切点。
刚体内相距 $u\uvec\tau\,\dd t$ 的两点，其速度相差
$\omega_i\uvec e_z\times(u\uvec\tau\,\dd t)=\omega_i u\uvec n_1\,\dd t$。
因此沿几何切点求导有
\[
 \frac{\dd}{\dd t}\vect v_i(P_3,t)=\vect a_i+\omega_i u\uvec n_1.
\]
对两个物体相减，利用接触材料点的速度持续相等，得到
\[
 \vect a_1-\vect a_2=(\omega_2-\omega_1)u\uvec n_1.
\]
分别沿两个内法向投影，记 $a_{1n}=\vect a_1\cdot\uvec n_1$、
$a_{2n}=\vect a_2\cdot\uvec n_2$，再代入 $u$，便有
\begin{equation}
  \boxed{a_{1n}+a_{2n}
  =\frac{(\omega_1-\omega_2)^2}{\kappa_1+\kappa_2}}.
  \label{r2ch01:eq:041}
\end{equation}
这里 $a_{1n},a_{2n}$ 分别沿指向各自物体内部的法向为正。右端非负，描述两个当前接触的材料点随后沿各自内部法向分开的相对趋势。若改用带符号曲率或统一公法线，公式的符号外观会改变，但从速度合成重新推导即可避免混淆。

以图~\ref{r2ch01:fig:010} 的半径 $R$ 圆盘为物体1、静止地面为物体2，取向右为切向正方向。
此时 $\kappa_1=1/R,\kappa_2=0,\omega_1=-v_O/R,\omega_2=0,\vect v_m=0$，
从而 $u=v_O$，$\vect v_{P_3}=v_O\uvec\tau$，并有 $a_{1n}=v_O^2/R$。
这同时检验了速度符号、加速度方向和曲率所在的位置。

当一个轮廓局部近似直线时，可以令该曲率为零；当两个轮廓的曲率都为零时，
上述分式成为退化式，不能除以 $\kappa_1+\kappa_2=0$。保持平行接触要求两者角速度相同，
接触位置的选择需由其余几何约束决定。若曲率和很小但不为零，相同相对角速度可使几何切点
相对轮廓快速移动；这不是材料点速度无限大的结论。

内接触可通过带符号曲率统一处理。保持前述 $\uvec\tau,\uvec n_1,\uvec n_2$ 的约定，
把曲线1沿 $\uvec\tau$ 前进时的切向变化写为 $k_1\uvec n_1$，曲线2的写为 $k_2\uvec n_2$。
曲面朝物体内部弯曲时 $k_i>0$，朝外部弯曲时 $k_i<0$。公切线转率仍满足
\[
 \dot\psi=\omega_1+k_1u=\omega_2-k_2u,
 \quad
 u=\frac{\omega_2-\omega_1}{k_1+k_2},
 \quad
 a_{1n}+a_{2n}=\frac{(\omega_1-\omega_2)^2}{k_1+k_2}.
\]
例如半径 $r$ 的小圆在半径 $R>r$ 的凹圆面内部纯滚时，$k_1=1/r$、$k_2=-1/R$，
分母变为 $1/r-1/R>0$，而不是 $1/r+1/R$。非退化的光滑接触通常要求该局部曲率差为正；
若 $r=R$，轮廓在这一级近似下相合，孤立几何切点不再由该公式确定，必须另行处理，不能声称分母必然严格大于零。

\subsection{刚杆约束}
\label{r2ch01:sec:rod}

长度为 $l$ 的刚杆连接端点 $1,2$。取沿杆单位矢量 $\uvec e$ 和垂直于杆的单位矢量 $\uvec e_\perp$，由 $\vect r_1-\vect r_2=l\uvec e$ 求导，得到
\begin{equation}
  v_{1\parallel}-v_{2\parallel}=0,
  \qquad
  v_{1\perp}-v_{2\perp}=\omega l.
  \label{r2ch01:eq:042}
\end{equation}
再求导并注意基矢也在转动，得到
\begin{equation}
  a_{1\parallel}-a_{2\parallel}=-\omega^2l,
  \qquad
  a_{1\perp}-a_{2\perp}=\beta l.
  \label{r2ch01:eq:043}
\end{equation}
速度约束中沿杆分量相等，并不意味着沿杆加速度分量相等；差值 $-\omega^2l$ 正是转动所需的向心项。

式~\eqref{r2ch01:eq:042}、\eqref{r2ch01:eq:043} 也可直接从长度约束
$|\vect r_1-\vect r_2|^2=l^2$ 连续求导得到。一次求导给
$(\vect v_1-\vect v_2)\cdot(\vect r_1-\vect r_2)=0$，说明相对速度垂直于杆；二次求导给
\[
  (\vect a_1-\vect a_2)\cdot(\vect r_1-\vect r_2)
  =-|\vect v_1-\vect v_2|^2,
\]
右端正是向心项的来源。这种做法不必预先画出角速度方向，适合三维杆约束。

\subsection{不可伸长柔绳约束}
\label{r2ch01:sec:rope}

设绳的物质坐标为 $s$，单位切向量为 $\uvec\tau$，速度场为 $\vect v(s,t)$。相邻绳元长度保持不变要求
\begin{equation}
  \frac{\partial\vect v}{\partial s}\cdot\uvec\tau=0.
  \label{r2ch01:eq:044}
\end{equation}
应特别注意，式~\eqref{r2ch01:eq:044} 并不是 $\partial(\vect v\cdot\uvec\tau)/\partial s=0$，因为后一式还包含 $\vect v\cdot\partial\uvec\tau/\partial s$。
在曲率非零处，令 $v_n=\vect v\cdot\uvec n$，乘积求导给出
\[
 \frac{\partial v_\tau}{\partial s}
 =\frac{\partial\vect v}{\partial s}\cdot\uvec\tau
   +\vect v\cdot\frac{\partial\uvec\tau}{\partial s}
 =\kappa v_n.
\]
因此切向速率沿绳处处相同还要求 $\kappa v_n=0$；直绳满足此条件，但一般正在改变形状的弯绳未必满足。
对由若干直段和动滑轮组成的系统，将各段长度写成坐标函数并对总长度求导，通常更稳妥。

这个局部条件可从相邻绳元的距离直接推导。设绳上相邻两个物质点相差 $\dd s$，分离矢量为
$\dd\vect r=\uvec\tau\,\dd s$。不可伸长要求 $|\dd\vect r|$ 不随时间改变，对 $|\dd\vect r|^2$ 求导得
\[
  0=2\dd\vect r\cdot\dd\vect v
  =2\uvec\tau\cdot\frac{\partial\vect v}{\partial s}(\dd s)^2,
\]
即式~\eqref{r2ch01:eq:044}。它只限制沿绳方向的相对伸缩率，允许相邻绳元的速度因绳的弯曲而有不同法向分量。所以弯绳可以改变形状，却不能沿局部切向被拉长。

若轻绳保持绷直且某一小段始终为直线，它在该段的运动学与不可伸长轻杆相同，端点沿绳方向速度相等。绳绕过固定滑轮时，接触弧长不随时间改变，滑轮两侧绳速大小相等；若滑轮本身移动，则必须把两侧可变直段和包角部分一起写入总长度。对于弯曲柔绳，局部切向会随 $s$ 改变，所以只能使用式~\eqref{r2ch01:eq:044}，不能把“切向速度分量处处相同”机械推广到任意曲线绳。

本章后续练习把上述工具用于抛体、追击、曲率、瞬时轴、复合滚动和机构约束。解答中凡使用活动标架或接触点，均明确说明所指的是物质点还是几何点。
